\documentclass[12pt,a4paper]{ctexart}

% ===== 页面设置 =====
\usepackage[top=2.5cm,bottom=2.5cm,left=2.5cm,right=2.5cm,headheight=15pt]{geometry}
\usepackage{amsmath,amssymb,amsthm}
\usepackage{enumitem}
\usepackage{xcolor}
\usepackage{tikz}
\usepackage{hyperref}
\usepackage{fancyhdr}

% ===== 页眉页脚 =====
\pagestyle{fancy}
\fancyhf{}
\fancyhead[L]{\textbf{高等数学（第七版·下册）——曲线积分的应用}}
\fancyhead[R]{\textbf{旋转曲面面积}}
\fancyfoot[C]{\thepage}
\renewcommand{\headrulewidth}{0.8pt}

% ===== 方框答案 =====
\newcommand{\answerbox}[1]{%
  \begin{center}
  \fbox{\color{red}\large\bfseries #1}
  \end{center}
}

% ===== 文档信息 =====
\title{\Huge\bfseries 曲线积分求旋转曲面面积}
\author{高等数学（同济大学数学系~编·第七版·下册）}
\date{\today}

\begin{document}

\maketitle
\thispagestyle{fancy}

% ==================== 原题 ====================
\section*{题目}
\fbox{%
  \parbox{\dimexpr\textwidth-2\fboxsep-2\fboxrule\relax}{%
    \textbf{3.}\; 用曲线积分求曲线段\;
    $y = x^2,\; 0 \leq x \leq 2$\;
    绕 $y$ 轴旋转所得的曲面面积。%
  }%
}
\bigskip

% ==================== 公式推导 ====================
\section*{第一步：如何把题目转化为 $\displaystyle\int_L \cdots ds$ 的积分？}

这是整个解题的\textbf{核心思想}。我们采用微元法，分四层理解。

\medskip

\noindent\textbf{\large ① 先把曲线切成无数小段}

将曲线 $L$ 用分点 $M_0,M_1,\ldots,M_n$ 分成 $n$ 个小弧段。取其中第 $i$ 小段
$\widehat{M_{i-1}M_i}$，记其弧长为 $\Delta s_i$。

\begin{center}
\begin{tikzpicture}[scale=1.8]
  % 坐标轴
  \draw[->,gray] (-0.3,0) -- (3.5,0) node[right] {$x$};
  \draw[->,gray] (0,-0.3) -- (0,2.5) node[above] {$y$};
  % 曲线
  \draw[thick,blue,samples=60,domain=0:1.7] plot (\x,{\x*\x});
  % 取一小段，放大
  \draw[very thick,red] (1.0,{1.0*1.0}) -- (1.25,{1.25*1.25});
  % 标注
  \filldraw (1.0,1.0) circle (0.03) node[below left] {$M_{i-1}$};
  \filldraw (1.25,1.5625) circle (0.03) node[right] {$M_i$};
  \draw[<->,red] (1.35,1.3) -- node[right] {$\Delta s_i$} (1.35,1.55);
  \node[blue] at (1.7,2.5) {$y=x^2$};
  \node at (2,0) [below] {$x$ 轴};
  % x_i
  \draw[dashed,gray] (1.0,0) -- (1.0,1.0);
  \draw[dashed,gray] (1.25,0) -- (1.25,1.5625);
  \node at (1.0,0) [below] {$x_{i-1}$};
  \node at (1.25,0) [below] {$x_i$};
\end{tikzpicture}
\end{center}

\medskip

\noindent\textbf{\large ② 每个小弧段绕 $y$ 轴旋转，形成什么？}

将第 $i$ 小段 $\widehat{M_{i-1}M_i}$ 绕 $y$ 轴旋转一整圈。

\emph{这段小弧很短}，可近似看作一条\textbf{倾斜的小线段}。它绕 $y$ 轴旋转后，
形成一个\textbf{窄圆环带}（想象一个截头圆锥的侧面）：

\begin{center}
\begin{tikzpicture}[scale=1.3]
  % Draw a perspective ring
  \draw[fill=blue!10] (0,2) ellipse (2 and 0.5);
  \draw[fill=blue!10] (0,2.5) ellipse (2.3 and 0.55);
  % sides
  \draw (2,2) -- (2.3,2.5);
  \draw (-2,2) -- (-2.3,2.5);
  % hidden lines
  \draw[dashed] (2.3,2.5) arc (0:180:2.3 and 0.55);
  \draw (1.8,3.5) -- (2.1,2.75);
  \draw[->] (2.3,2.5) arc (0:-30:2.3 and 0.55) node[midway,right] {旋转};
  % y-axis
  \draw[->,gray,thick] (0,0.5) -- (0,4) node[left] {$y$};
  % labels
  \node at (1.8,2.25) {$2\pi x_i$ (周长)};
  \draw[<->,red] (-2.3,2.8) -- node[above] {$\Delta s_i$ (弧宽)} (2.3,2.8);
\end{tikzpicture}
\end{center}

\textbf{关键问题：这个窄环带的面积是多少？}

\newpage

\noindent\textbf{\large ③ 近似计算窄环带的面积}

因为 $\Delta s_i$ 很短，环带上下几乎一样大。把它沿着母线剪开摊平，
近似为一个\textbf{矩形}：

\[
\Delta S_i \;\approx\; \underbrace{(2\pi \times \text{半径})}_{\text{周长}}\, \times \,
                         \underbrace{\Delta s_i}_{\text{宽度}}
               \;=\; 2\pi \cdot \bar{x}_i \cdot \Delta s_i
\]

其中 $\bar{x}_i$ 是第 $i$ 小段上某点的 $x$ 坐标（可取该段上任一点的 $x$）。

\begin{center}
\fbox{%
\begin{minipage}{0.7\textwidth}
\centering
\small
$\Delta S_i$ = 环带侧面积 $\approx$（周长）$\times$（宽度）
= $2\pi x_i \cdot \Delta s_i$
\end{minipage}%
}
\end{center}

\medskip

\noindent\textbf{\large ④ 求和，取极限 $\to$ 积分}

将每一小段的面积累加：
\[
S \approx \sum_{i=1}^n 2\pi \bar{x}_i \,\Delta s_i
\]

令分点无限加密（$n\to\infty$，$\Delta s_i \to 0$），上式取极限即得\textbf{第一类曲线积分}：
\[
\boxed{S = 2\pi \int_L x \, ds}
\]

这就是公式的由来！其中 $ds$ 是弧长微元，
对函数 $y=f(x)$ 有 $ds = \sqrt{1+[f'(x)]^2}\,dx$。

% ==================== 两种旋转轴对照 ====================
\subsection*{绕不同轴的公式对比}

\begin{center}
\renewcommand{\arraystretch}{1.8}
\begin{tabular}{|c|c|c|}
  \hline
  \textbf{旋转轴} & \textbf{曲面面积公式} & \textbf{直观理解} \\
  \hline
  绕 $y$ 轴 & $S = 2\pi \displaystyle\int_L x\, ds$ & 半径 $= x$，周长 $=2\pi x$ \\
  \hline
  绕 $x$ 轴 & $S = 2\pi \displaystyle\int_L |y|\, ds$ & 半径 $= |y|$，周长 $=2\pi|y|$ \\
  \hline
\end{tabular}
\end{center}

\textbf{记忆方法：} 面积 = 曲线上的点离旋转轴的\textbf{距离}$\times 2\pi\times$弧长微元 再积分。

简单说：\textcolor{red}{绕哪个轴转，就乘那点到该轴的距离的 $2\pi$ 倍}。

% ==================== 解答 ====================
\section*{解答}

\noindent\textbf{步骤1：求 $y'(x)$ 并计算弧长微元 $ds$}
\[
y = x^2 \quad\Longrightarrow\quad y'(x) = 2x.
\]
\[
ds = \sqrt{1 + \bigl[y'(x)\bigr]^2}\, dx
    = \sqrt{1 + (2x)^2}\, dx
    = \sqrt{1 + 4x^2}\, dx.
\]

\medskip

\noindent\textbf{步骤2：代入旋转曲面面积公式}
\[
S = 2\pi \int_L x\, ds
   = 2\pi \int_0^2 x \cdot \sqrt{1 + 4x^2}\, dx.
\]

\medskip

\noindent\textbf{步骤3：换元积分}

令 $u = 1 + 4x^2$，则：
\[
du = 8x\, dx \quad\Longrightarrow\quad x\, dx = \frac{1}{8}\, du.
\]

积分限变换：
\[
x = 0 \;\to\; u = 1,\qquad
x = 2 \;\to\; u = 1 + 4 \times 4 = 17.
\]

代入：
\[
\begin{aligned}
S &= 2\pi \int_1^{17} \sqrt{u} \cdot \frac{1}{8}\, du
    = \frac{\pi}{4} \int_1^{17} u^{\frac{1}{2}}\, du \\[6pt]
  &= \frac{\pi}{4} \cdot \frac{2}{3}\, u^{\frac{3}{2}} \Biggr|_1^{17} \\[6pt]
  &= \frac{\pi}{6} \Bigl( 17^{\frac{3}{2}} - 1^{\frac{3}{2}} \Bigr).
\end{aligned}
\]

\medskip

\noindent\textbf{步骤4：化简结果}
\[
17^{\frac{3}{2}} = 17 \cdot \sqrt{17},
\]
\[
\therefore\;
S = \frac{\pi}{6} \bigl( 17\sqrt{17} - 1 \bigr).
\]

\bigskip

\answerbox{$\displaystyle S = \frac{\pi}{6}\bigl(17\sqrt{17} - 1\bigr)$}

\bigskip

% ==================== 数值验证 ====================
\section*{数值结果}

\[
\sqrt{17} \approx 4.1231,\quad
17\sqrt{17} \approx 70.09,
\]
\[
S \approx \frac{\pi}{6} \times 69.09 \approx 36.18.
\]

即 $S \approx 36.2$（平方单位）。

% ==================== 方法小结 ====================
\section*{方法小结}

\begin{enumerate}[leftmargin=*,itemsep=4pt]
  \item \textbf{识别公式：} 绕 $y$ 轴旋转 $\to S = 2\pi\int_L x\,ds$\;
        （绕 $x$ 轴则是 $S = 2\pi\int_L |y|\,ds$）
  \item \textbf{求 $ds$：} $ds = \sqrt{1+(y')^2}\,dx$
  \item \textbf{换元积分：} 凑微分 $x\,dx = \frac{1}{8}du$，化为 $\int \sqrt{u}\,du$ 型
\end{enumerate}

此题属于曲线积分的几何应用，在教材\S11.1 中与弧长、质心、转动惯量等内容并列，
本质上是将平面曲线的弧长积分与旋转体的侧面积公式结合。

\end{document}
